\chapter{Implementação}
\label{chapter:implementacao}

O objetivo deste capítulo é descrever como a Plataforma de Engenharia Self-Service — ProxLab — foi implementada, detalhando a arquitetura do \textit{backend}, os fluxos de provisionamento de máquinas virtuais e de contêineres, e os mecanismos de autenticação e acesso remoto oferecidos ao usuário final, conforme a visão geral da Figura~\ref{fig:arquitetura_proxlab}.

\figimg{fig:arquitetura_proxlab}{Arquitetura geral do ProxLab}{arquitetura_proxlab}{elaboração própria}

\section{Arquitetura em camadas}

O \textit{backend} foi desenvolvido em Go, utilizando o \textit{framework} web Gin e o ORM GORM sobre um banco de dados PostgreSQL. A aplicação segue uma arquitetura em camadas, organizada de forma que cada camada dependa apenas de interfaces, e não de implementações concretas (Figura~\ref{fig:camadas_backend}):

\begin{itemize}
    \item \texttt{handler} — controladores HTTP (Gin), responsáveis por validar o \textit{payload} da requisição e delegar a lógica de negócio;
    \item \texttt{usecase} — lógica de negócio da aplicação, independente de detalhes de infraestrutura;
    \item \texttt{domain/repository} — interfaces de persistência, nunca implementações;
    \item \texttt{repository} — implementações concretas das interfaces de persistência, utilizando GORM.
\end{itemize}

Essa separação é reforçada por uma regra de desenvolvimento adotada no projeto: nenhuma lógica de negócio deve residir no ponto de entrada da aplicação, e nenhuma chamada ao GORM deve ocorrer diretamente em um \textit{handler}.

\figimg{fig:camadas_backend}{Arquitetura em camadas do \textit{backend}}{camadas_backend}{elaboração própria}

\section{Provisionamento de máquinas virtuais}

O fluxo de criação de uma máquina virtual segue o caminho \textit{Handler} \(\rightarrow\) \textit{UseCase} \(\rightarrow\) \textit{Provisioner} \(\rightarrow\) \texttt{terraform-exec} \(\rightarrow\) Proxmox VE. Ao receber uma solicitação de criação de VM, o \textit{backend} gera um arquivo de configuração Terraform (\texttt{main.tf}) em um \textit{workspace} isolado por identificador único de VM, utilizando o \textit{provider} \texttt{bpg/proxmox}, e executa o Terraform por meio da biblioteca \texttt{terraform-exec}. As credenciais de acesso ao Proxmox são fornecidas apenas via variáveis de ambiente do processo, nunca gravadas em arquivo de configuração.\\

\par Máquinas virtuais Ubuntu são criadas a partir da clonagem completa de um \textit{template} \textit{cloud-init} previamente preparado no Proxmox, já nascendo ligadas e com um usuário padrão cuja senha é gerada de forma aleatória por VM e devolvida uma única vez na resposta da criação, nunca persistida em banco de dados. Máquinas virtuais de outros sistemas operacionais seguem o fluxo de VM vazia, cabendo ao usuário realizar a instalação do sistema operacional.\\

\par Quando o usuário solicita GPU, o \textit{backend} escolhe um \textit{resource mapping} PCI do Proxmox VE (o que tiver menos máquinas associadas) e o anexa à VM por \textit{pass-through}, por meio de um bloco \texttt{hostpci} na configuração Terraform. Como a placa não pode ser compartilhada, a VM com GPU nasce desligada e a operação de ligar é serializada: antes de iniciá-la, o \textit{backend} verifica se outra VM está usando o mesmo \textit{mapping} e, se o Proxmox VE recusar a inicialização, a falha é traduzida em uma resposta de conflito, sem expor ao usuário a mensagem interna da infraestrutura. A instalação automática do \textit{driver} de vídeo e do Docker nessas VMs, por \textit{cloud-init}, é descrita no Apêndice~\ref{apendice:material}.\\

\par Operações de consulta de status, ligar e desligar uma VM já existente não passam pelo Terraform: são realizadas diretamente contra a API REST do Proxmox VE, por meio de um cliente HTTP dedicado (\textit{ProxmoxClient}), autenticado via \textit{token} de API. A Figura~\ref{fig:tela_criar_vm} apresenta a tela de criação de máquina virtual utilizando o portal do ProxLab.

\figimg{fig:tela_criar_vm}{Tela de criação de máquina virtual no portal}{front-tela-criar-vm}{captura de tela do ProxLab}

\section{Provisionamento de contêineres de aplicação}

Contêineres de aplicação seguem um fluxo distinto: \textit{Handler} \(\rightarrow\) \textit{UseCase} \(\rightarrow\) \textit{ContainerRuntime} \(\rightarrow\) \texttt{client-go} \(\rightarrow\) cluster k3s. O cluster é executado em uma máquina virtual dedicada, provisionada uma única vez por meio de um utilitário de \textit{bootstrap} que instala o k3s e gera o \textit{kubeconfig} utilizado pelo \textit{backend}.\\

\par Cada usuário da plataforma possui um \textit{namespace} próprio no cluster, identificado por um sufixo derivado do seu identificador, e uma quota de recursos (CPU, memória) definida de acordo com o seu perfil (aluno ou pesquisador). O estado de um contêiner não é persistido em banco de dados: ele é lido diretamente do cluster a cada requisição, garantindo que a plataforma sempre reflita o estado real do recurso.\\

\par A exposição de rede de um contêiner depende do tipo de acesso solicitado pelo usuário: quando o contêiner expõe uma porta HTTP e um domínio de \textit{Ingress} está configurado, o acesso é feito por meio de um subdomínio dedicado; quando o acesso é via TCP, ou quando não há domínio configurado, o serviço é exposto por meio de \textit{NodePort}, em uma porta na faixa 30000–32767 do nó do cluster. Um contêiner também pode ser criado apenas para acesso via terminal, sem nenhuma porta de rede exposta. A Figura~\ref{fig:tela_criar_container} apresenta a tela de criação de contêineres usando o portal do ProxLab.

\figimg{fig:tela_criar_container}{Tela de criação de contêiner de aplicação no portal}{front-tela-criar-container}{captura de tela do ProxLab}

\section{Autenticação e gestão de usuários}

O acesso à plataforma é protegido por autenticação baseada em JWT. Após o \textit{login} (Figura~\ref{fig:tela_login}), o token é armazenado no \texttt{localStorage} do navegador e enviado em todas as requisições subsequentes por meio do cabeçalho \texttt{Authorization}. Conexões WebSocket, por não suportarem cabeçalhos HTTP customizados na negociação inicial, recebem o token por meio de um parâmetro de consulta (\textit{query param}) na URL de conexão. Não há autocadastro de usuários: o cadastro de novos usuários exige um token JWT de um usuário com perfil \textit{pesquisador}. O perfil define os limites de recursos e o \textit{pool} do Proxmox em que as VMs são criadas, sempre decidido pelo \textit{backend} e nunca pelo \textit{payload} da requisição. Os usuários podem alterar seu nome e senha em uma tela de perfil, e o pesquisador dispõe de uma tela de gestão que lista, filtra por perfil, altera o e-mail e exclui usuários (Figura~\ref{fig:tela_usuarios}); a exclusão remove antes todas as VMs e contêineres do usuário.

\figimg{fig:tela_login}{Tela de \textit{login} da plataforma}{front-tela-login(geral)}{captura de tela do ProxLab}

\figimg{fig:tela_usuarios}{Tela de usuários cadastrados (visão do pesquisador)}{front-tela-usuarios(pesquisador)}{captura de tela do ProxLab}

\section{Terminal web}

O acesso interativo aos recursos provisionados é oferecido diretamente pelo navegador, sem a necessidade de um cliente SSH local (Figura~\ref{fig:terminal_web}). Para máquinas virtuais, o \textit{backend} atua como um proxy WebSocket entre o navegador e o protocolo \textit{termproxy} do Proxmox VE, reproduzindo o seu \textit{handshake} específico: uma primeira mensagem de autenticação no formato \texttt{"<usuário>:<ticket>\textbackslash n"}, respondida com \texttt{OK}, seguida do enquadramento de dados de teclado e de redimensionamento de terminal segundo o protocolo do Proxmox. Para contêineres, o acesso a terminal é feito por meio da operação de \textit{exec} do próprio cluster Kubernetes, sem a necessidade de um protocolo intermediário.

\figimgh{fig:terminal_web}{Terminal web de uma máquina virtual Ubuntu no navegador}{front-tela-terminal-web-vm}{captura de tela do ProxLab}{0.17}

\section{Infraestrutura como código e configuração inicial}

A infraestrutura de máquinas virtuais é definida como código: cada VM possui seu próprio \textit{workspace} Terraform, de modo que criar ou destruir um recurso equivale a aplicar ou destruir uma configuração versionável e reprodutível. A configuração inicial das VMs Ubuntu (usuário, senha e chaves SSH) é entregue por \textit{cloud-init}, a partir do \textit{template} clonado. Não é utilizado o Ansible: a instalação de \textit{software} é oferecida por meio de contêineres de aplicação. A Figura~\ref{fig:terraform_main_tf} apresenta um exemplo de configuração usando o Terraform.

\figimgh{fig:terraform_main_tf}{Exemplo de configuração Terraform gerada pelo \textit{backend}}{terraform_main_tf}{elaboração própria (exemplo ilustrativo baseado no \textit{template} do \textit{backend})}{0.16}
